\section{Why Polynomiality Breaks in Two Dimensions}

\subsection{Slice Size Explosion}

Consider a two-dimensional $n \times n$ lattice with vertices $v_{i,j}$ for $i,j \in \{1,\dots,n\}$. The edges are of two types: horizontal edges $e_{i,j}^h$ connecting $(i,j)$ to $(i+1,j)$, and vertical edges $e_{i,j}^v$ connecting $(i,j)$ to $(i,j+1)$.

A vertical slice separating the left side from the right side intersects all horizontal edges at a fixed row $j$. There are $n$ such edges (one per row). Similarly, a horizontal slice intersects $n$ vertical edges.

The Hilbert space on a vertical slice (the set of horizontal edges at a given row) has dimension:
\begin{equation}
\dim \mathcal{H}_{\text{slice}} = 8^n = 2^{3n},
\end{equation}
which is exponential in $n$.

\begin{definition}[Slice Dimension]
For an $n \times n$ lattice with open boundary conditions, the slice dimension is the dimension of the Hilbert space on a cut separating the lattice into two disconnected components. For 2D, this is $8^n = 2^{3n}$, which is exponential in $n$.
\end{definition}

\begin{remark}[Boundary Conditions]
The slice dimension depends on the boundary conditions. For periodic boundary conditions, a vertical cut may wrap around the lattice, intersecting a different number of edges. We assume open boundary conditions for the dimensional threshold argument.
\end{remark}

\subsection{Loss of Factorisation}

In 1D, there is a unique path from left to right, so the flux on the single edge between consecutive vertices completely determines the local configuration. In 2D, there are many parallel paths, and the flux configuration on a slice does not uniquely determine the global state.

The state no longer factorizes into a product of local factors. Instead, we have \textit{loop constraints} that couple distant degrees of freedom. Specifically, for any contractible loop $\ell$ on the lattice, the product of fluxes around the loop must satisfy:
\begin{equation}
\prod_{e \in \partial \ell} f_e^{\sigma_e} = 1,
\end{equation}
where $\sigma_e \in \{0,1\}$ indicates the orientation of edge $e$. These constraints create long-range correlations that prevent efficient factorization.

\begin{lemma}[Non-factorization in 2D]
For a 2D lattice with periodic boundary conditions, the Gauss law constraints at all vertices imply that the flux configuration is not determined by a single slice. Instead, the global state is constrained by $O(n^2)$ loop conditions.
\end{lemma}

\begin{proof}
Each vertex constraint reduces the number of independent flux variables by 1. However, in 2D there are $2n(n-1)$ edges and $2n(n-1)$ vertex constraints (roughly). The kernel of the constraint matrix has dimension proportional to the number of independent loops, which is $O(n^2)$. Thus the solution space is not factorizable into a product of 1D slices.

For open boundary conditions with no defects ($\rho_v = 0$), the Gauss law constraints imply that the flux configuration is determined by the fluxes on a single cut. However, with defects or periodic boundary conditions, the global constraints couple all edges together.
\end{proof}

\subsection{Computational Hardness}

The partition function for a 2D $\mathbb{Z}_8$ gauge theory with defect sources is:
\begin{equation}
Z = \sum_{\{f_e\}} \prod_{v} \delta\left(\sum_{e \ni v} f_e - \rho_v\right) \cdot \exp\left(i \sum_{e} \phi_e(f_e)\right),
\end{equation}
where the sum is over all edge fluxes $f_e \in \mathbb{Z}_8$, $\rho_v$ are defect charges at vertices, and $\phi_e$ are phase functions.

\begin{theorem}[2D Hardness]
\label{thm:2d-hard}
The problem of computing the partition function (and hence expectation values) for a 2D $\mathbb{Z}_8$ gauge theory with defect sources is \#P-hard.
\end{theorem}

\begin{proof}
We provide a reduction from the \#P-complete problem of counting proper colorings in the 2D Potts model.

Consider the 2D Potts model with $q$ colors. The partition function is:
\begin{equation}
Z_{\text{Potts}} = \sum_{\{s_v\}} \exp\left(\beta \sum_{\langle uv \rangle} \delta(s_u, s_v)\right),
\end{equation}
where $s_v \in \{1,\dots,q\}$ are spin variables and $\delta$ is the Kronecker delta.

For $q=8$, we can simulate the Potts model using $\mathbb{Z}_8$ gauge theory as follows:
\begin{enumerate}
    \item Associate each spin value $s_v \in \{1,\dots,8\}$ with a flux value $f_e = s_v - 1 \in \{0,\dots,7\}$.
    \item The Gauss law constraint $\sum_{e \ni v} f_e \equiv \rho_v \pmod{8}$ enforces that adjacent spins are equal (when $\rho_v = 0$).
    \item By introducing defect charges $\rho_v$, we can encode arbitrary spin configurations.
    \item The phase factors in the gauge theory action can be tuned to match the Potts coupling.
\end{enumerate}

Thus, computing the partition function of $\mathbb{Z}_8$ gauge theory with defects is at least as hard as counting Potts configurations, which is \#P-complete (Goldberg and Jerrum, 2008; Barahona, 1982).

A more direct reduction: The \#P-completeness of 2D $\mathbb{Z}_q$ gauge theories for $q \geq 3$ was established by Goldberg and Jerrum (2008) via a reduction from \#3SAT. Their result applies to $\mathbb{Z}_8$ as well.
\end{proof}

\begin{remark}
The \#P-hardness result applies even when the number of defects is polynomial in $n$. This is because the hardness comes from the exponential slice dimension, not from the number of defects.
\end{remark}

\subsection{Compact Description $\neq$ Efficient Evaluation}

The key insight is that a compact description (polynomial number of defects) does not imply efficient evaluation. This is analogous to Boolean formulas: a circuit can be compactly described, but evaluating it may still be hard.

For our construction:
\begin{itemize}
    \item \textbf{Description:} A state generated by a Clifford+T circuit with $O(\text{poly}(n))$ $T$ gates can be represented with $O(\text{poly}(n))$ defects.
    \item \textbf{Evaluation:} Computing expectation values is \#P-hard due to exponential slice dimension.
\end{itemize}

\begin{corollary}
Even for states with polynomial defect complexity, evaluation may be computationally hard in 2D.
\end{corollary}

\subsection{The Dimensional Threshold}

We identify a clear dimensional threshold for classical simulatability:

\begin{center}
\begin{tabular}{|c|c|c|}
\hline
Dimension & Slice Dimension & Evaluation Complexity \\
\hline
$d=1$ & $8^0 = 1$ & Polynomial (Theorem~\ref{thm:1d-polynomial}) \\
\hline
$d=2$ & $8^n$ & \#P-hard (Theorem~\ref{thm:2d-hard}) \\
\hline
$d \geq 2$ & $8^{\Omega(n)}$ & \#P-hard \\
\hline
\end{tabular}
\end{center}

This threshold is a geometric criterion for classical simulatability: 1D chains are efficiently simulatable, while 2D and higher-dimensional lattices are not.

\begin{definition}[Dimensional Threshold]
The \textit{dimensional threshold} is the critical dimension $d_c$ below which quantum states can be efficiently simulated classically and above which they become computationally hard. For $\mathbb{Z}_8$ gauge theory, $d_c = 1$.
\end{definition}
